
La construcción de circuitos biohíbridos requiere el uso de software con restricciones temporales que garanticen que la actividad registrada de la neurona viva y la implementación del modelo de neurona artificial se realizan con una tasa de transferencia de al menos $10$ $kHz$. En la mayor parte de los casos, esto significa que hay que utilizar herramientas de tiempo real estricto.

Durante este trabajo, se han descrito diversos sistemas y desarrollado varias herramientas en tiempo real para crear y estudiar la dinámica neuronal. De estos sistemas, se ha puesto especial atención en \textit{Preempt-RT} como \textit{RTOS}. Esta decisión viene dada, en parte por el soporte, fiabilidad y seguridad que la \textit{Linux Fundation} aporta a dicho sistema, pero también porque permite usar herramientas actualizadas que pueden estar disponibles únicamente en las últimas versiones de Linux; por otro lado, respecto a las herramientas, se ha centrado el trabajo en \textit{RTXI}, dicho enfoque es debido a su utilidad a la hora de realizar experimentos con circuitos biohíbridos, implementando no solo los módulos necesarios para ser usados en el nuevo kernel \textit{realtime} de Linux, sino también migrando aquellas herramientas de \textit{RTHybrid} que estaban en la antigua versión del programa, para que puedan ser usadas por los investigadores del \textit{GNB}.

Estas implementaciones han sido validadas de dos formas: La primera respecto al rendimiento del sistema, que ha sido probado con las herramientas de medición de \textit{RTXI}, validando su correcto funcionamiento y cumplimiento con las restricciones temporales. La segunda mediante el programa \textit{cyclictest}, que muestra la latencia del sistema en sí. Las latencias presentadas en este último también son admisibles.

Otra forma adicional de validación se ha realizado mediante la creación de circuitos "biohíbridos". Si bien los circuitos presentados no tienen un componente biológico real, sí que demuestran la posibilidad de crear un circuito biohíbrido. El primer circuito se compone de la parte virtual (modelo de Hindmarsh-Rose en \textit{RTXI}) y una neurona electrónica que funciona a la misma velocidad que una neurona viva. Su interacción mantenida el tiempo demuestra que el sistema puede seguir el ritmo real de un agente externo. El otro tipo de circuito es la interacción de la neurona virtual con la grabación de una neurona viva, conectando el modelo de \textit{RTXI} de manera unidireccional con dicha grabación.

Además, la capacidad de RTXI para operar con modelos neuronales complejos en tiempo real, como el de Hindmarsh-Rose, y su adaptación al kernel \textit{Preempt-RT}, posiciona esta plataforma como una solución para futuros desarrollos en neurotecnología. La integración de estos sistemas no solo simplifica la interacción con células neuronales, sino que también abre la puerta a nuevas configuraciones experimentales, incluyendo la posibilidad de controlar dinámicamente sinapsis biohíbridas en función del estado de red, o aplicar algoritmos de aprendizaje en bucle cerrado directamente sobre el hardware experimental. Esta versatilidad amplía el abanico de aplicaciones en las que estas herramientas podrían ser fundamentales, desde el estudio de la plasticidad sináptica hasta el desarrollo de prótesis neuronales adaptativas.

Toda la metodología y software desarrollado están disponibles en el laboratorio del \textit{GNB} para implementar una gran variedad de circuitos biohíbridos con neuronas vivas. Ya están programados experimentos que utilizarán el trabajo desarrollado.

Si bien todavía hay margen de mejora, tanto por parte del kernel, como de \textit{RTXI} y sus correspondientes plugins, actualmente se encuentra en un estado funcional que permite su uso al menos de manera experimental. No se debe descartar la posibilidad de que, en un futuro cercano, estas opciones destaquen a la hora de investigar y crear tecnologías que fusionen aún más el mundo computacional y el neuronal.
